% !TeX program = xelatex
% Compile with XeLaTeX in Overleaf.

\documentclass[a4paper]{article}
\usepackage{xeCJK}
\setCJKmainfont[BoldFont=FandolSong-Bold.otf]{FandolSong-Regular.otf}
\usepackage{needspace}

% \usepackage[default]{fontsetup}

\usepackage{fancyhdr}
\usepackage{extramarks}
\usepackage{amsmath}
\usepackage{amsthm}
\usepackage{amsfonts}
\usepackage{tikz}
\usepackage[plain]{algorithm}
\usepackage{algpseudocode}
\usepackage{enumerate}
\usepackage{tikz}

\usetikzlibrary{automata,positioning}

%
% Basic Document Settings
%

\topmargin=-0.2in
\evensidemargin=0in
\oddsidemargin=0in
\textwidth=6.5in
\textheight=9.5in
\headsep=0.25in

\linespread{1.1}

\pagestyle{fancy}
\lhead{}
\chead{\hmwkClass : \hmwkTitle}
\rhead{\firstxmark}
\lfoot{\lastxmark}
\cfoot{\thepage}

\renewcommand\headrulewidth{0.4pt}
\renewcommand\footrulewidth{0.4pt}

\setlength\parindent{0pt}

%
% Create Problem Sections
%

\newcommand{\exitProblemHeader}[1]{
    \nobreak\extramarks{Problem \arabic{#1} }{Problem \arabic{#1} continued on next page\ldots}\nobreak{}
    \stepcounter{#1}
    \nobreak\extramarks{Problem \arabic{#1}}{}\nobreak{}
}

\newcommand*\circled[1]{\tikz[baseline=(char.base)]{
    \node[shape=circle,draw,inner sep=2pt] (char) {#1};}}

\setcounter{secnumdepth}{0}
\newcounter{partCounter}
\newcounter{homeworkProblemCounter}
\setcounter{homeworkProblemCounter}{1}
\nobreak\extramarks{Problem \arabic{homeworkProblemCounter}}{}\nobreak{}

%
% Homework Problem Environment
%

\newenvironment{homeworkProblem}[1][-1]{
    \ifnum#1>0
        \setcounter{homeworkProblemCounter}{#1}
    \fi
    \section{Problem \arabic{homeworkProblemCounter}}
    \setcounter{partCounter}{1}
}{
    \exitProblemHeader{homeworkProblemCounter}
}

%
% Homework Details
%

\newcommand{\hmwkTitle}{Homework\ \#05}
\newcommand{\hmwkClass}{Probability \& Statistics for EECS}
\newcommand{\hmwkDueDate}{2026-11-3}
\newcommand{\hmwkAuthorName}{杨涵}
\newcommand{\hmwkAuthorID}{2025531064}

%
% Title Page
%

\title{
    \vspace{2in}
    \textmd{\textbf{\hmwkClass:\\ \hmwkTitle}}\\
    \normalsize\vspace{0.1in}\small{Due\ on\ \hmwkDueDate\ at 18:00}\\
    \vspace{4in}
}

\author{
    Name: \textbf{\hmwkAuthorName} \\
    Student ID: \hmwkAuthorID}
\date{}

\renewcommand{\part}[1]{\textbf{\large Part \Alph{partCounter}}\stepcounter{partCounter}\\}

%
% Various Helper Commands
%

\newcommand{\alg}[1]{\textsc{\bfseries \footnotesize #1}}
\newcommand{\deriv}[1]{\frac{\mathrm{d}}{\mathrm{d}x} (#1)}
\newcommand{\pderiv}[2]{\frac{\partial}{\partial #1} (#2)}
\newcommand{\dx}{\mathrm{d}x}
\newcommand{\solution}{\par\addvspace{\baselineskip}\Needspace{4\baselineskip}\noindent\textbf{\large Solution}\par\nobreak\smallskip}
\newcommand{\E}{\mathrm{E}}
\newcommand{\Var}{\mathrm{Var}}
\newcommand{\Cov}{\mathrm{Cov}}
\newcommand{\Bias}{\mathrm{Bias}}

\begin{document}

\date{
    Due on Nov. 3, 2026, 18:00 UTC+8}
\title{SI 140A-01  Probability \& Statistics for EECS, Fall 2026 \\
    Homework 5}
\maketitle
Read all the instructions below carefully before you start working on the assignment, and before you make a submission.
\begin{itemize}
    \item You are required to write down all the major steps towards making your conclusions; otherwise you may obtain limited points of the problem.
    \item Write your homework in English; otherwise you will get no points of this homework.
    \item Any form of plagiarism will lead to $0$ point of this homework.
\end{itemize}

\newpage

\begin{homeworkProblem}
If $X \sim \text{HGeom}(w,b,n)$, what is the distribution of $n-X$? Give a short proof.

\solution
Since $X\sim\mathrm{HGeom}(w,b,n)$,
\[
P(X=k)=\frac{\binom{w}{k}\binom{b}{n-k}}{\binom{w+b}{n}}.
\]
Let $Y=n-X$. Then $X=n-Y$, so
\[
P(Y=k)=P(X=n-k)
=\frac{\binom{w}{n-k}\binom{b}{k}}{\binom{w+b}{n}},
\qquad \max(0,n-w)\leq k\leq\min(n,b).
\]
Therefore,
\[
n-X\sim\mathrm{HGeom}(b,w,n).
\]

\end{homeworkProblem}

\newpage

\begin{homeworkProblem}
A particular Youtuber is known for his arbitrary steak-eating habit during live streaming.
In each live streaming day, he orders a steak with doneness as one of well done, medium well, medium, medium rare, and rare (with equal probability, independent of other live streaming days).
How many days of live streaming do you expect to watch before you see him eating steaks with each possible doneness at least one once (suppose you are a big fan who watches every live streaming)?

\solution
After seeing $i-1$ distinct levels of doneness, let $W_i$ be the number of
additional days needed to see a new level, including the day on which it
appears. The total number of days is
\[
T=W_1+W_2+W_3+W_4+W_5.
\]

For independent trials with success probability $p$, let $W$ count the
trials up to and including the first success. Then
\[
P(W=t)=(1-p)^{t-1}p,\qquad t=1,2,3,\dots.
\]
To find its expectation, write $q=1-p$. Differentiating the geometric series
gives
\[
\sum_{t=0}^{\infty}q^t=\frac1{1-q}
\quad\Longrightarrow\quad
\sum_{t=1}^{\infty}tq^{t-1}=\frac1{(1-q)^2},
\qquad 0\leq q<1.
\]
Hence,
\[
\E[W]=\sum_{t=1}^{\infty}t(1-p)^{t-1}p
=\frac{p}{(1-q)^2}=\frac1p.
\]

For $W_i$, the probability of seeing a new level is
\[
p_i=\frac{5-(i-1)}5=\frac{6-i}5,
\qquad \E[W_i]=\frac5{6-i},\qquad i=1,\dots,5.
\]
By linearity of expectation,
\begin{align*}
\E[T]
&=\E[W_1]+\E[W_2]+\E[W_3]+\E[W_4]+\E[W_5]\\
&=1+\frac54+\frac53+\frac52+5\\
&=6+\frac{15+20+30}{12}
=\frac{137}{12}\approx11.42\text{ days}.
\end{align*}

\end{homeworkProblem}

\newpage

\begin{homeworkProblem}
Mario and Zelda are independently performing independent Bernoulli trials.
For concreteness, assume that Mario is flipping a nickel with probability $p_{1}$ of Heads and Zelda is flipping a penny with probability $p_2$ of Heads.
Let $X_{1}, X_{2}, \cdots $ be Mario's results and $Y_{1}, Y_{2}, \cdots $ be Zelda's results, with $X_{i} \sim \text{Bern}(p_{1})$ and $Y_{j} \sim \text{Bern}(p_{2})$.

\begin{enumerate}[(a)]
    \item Find the distribution and expected value of the first time at which they are simultaneously successful, \textit{i.e.}, the smallest $n$ such that $X_{n} = Y_{n} = 1$.
    \item Find the expected time until at least one has a success (including the success).
    \item For $p_{1} = p_{2}$, find the probability that their first successes are simultaneous, and use this to find the probability that Mario's first success precedes Zelda's.
\end{enumerate}

\solution
\textbf{(a)} Let $T$ be the first trial on which Mario and Zelda are
simultaneously successful. By independence, the probability of simultaneous
success in each trial is $p_1p_2$. Thus,
\[
P(T=n)=(1-p_1p_2)^{n-1}p_1p_2,
\qquad n=1,2,3,\dots.
\]
So $T$ has a geometric distribution on $\{1,2,3,\dots\}$ with success
probability $p_1p_2$. Using the expectation derived in Problem 2,
\[
\E[T]=\frac1{p_1p_2}.
\]

\medskip
\Needspace{9\baselineskip}
\textbf{(b)} The probability that both fail in a trial is
\[
q=(1-p_1)(1-p_2).
\]
Let $S$ be the first trial on which at least one succeeds. Its success
probability is $1-q$, so
\[
\E[S]=\frac1{1-q}=\frac1{p_1+p_2-p_1p_2}.
\]

\medskip
\Needspace{17\baselineskip}
\textbf{(c)} Let $p_1=p_2=p$, and let $M$ be the event that their first
successes are simultaneous. For this to happen on trial $n$, both must fail
in the first $n-1$ trials and both succeed on trial $n$. Therefore,
\begin{align*}
P(M)
&=\sum_{n=1}^{\infty}\bigl[(1-p)^2\bigr]^{n-1}p^2\\
&=p^2\bigl[1+(1-p)^2+(1-p)^4+\cdots\bigr]\\
&=\frac{p^2}{1-(1-p)^2}
=\frac{p^2}{p(2-p)}=\frac{p}{2-p}.
\end{align*}
By symmetry, Mario and Zelda are equally likely to have the earlier first
success. Let $p'$ be the probability that Mario's first success precedes
Zelda's. Then
\[
2p'+P(M)=1,
\qquad
p'=\frac12\left(1-\frac{p}{2-p}\right)=\frac{1-p}{2-p}.
\]

\end{homeworkProblem}

\newpage

\begin{homeworkProblem}
Show that for any events $A_1,\dots,A_n$,
\[
P(A_1 \cap A_2 \cap \cdots \cap A_n) \ge \sum_{j=1}^n P(A_j) - n + 1.
\]

Hint: First prove a similar statement about indicator r.v.s, by interpreting what the events $I(A_1 \cap A_2 \cdots \cap A_n)=1$ and $I(A_1 \cap A_2 \cdots \cap A_n)=0$ mean.

\solution
For any event $A$, define its indicator by
\[
I(A)=
\begin{cases}
1, & \text{if }A\text{ occurs},\\
0, & \text{if }A\text{ does not occur}.
\end{cases}
\]
First, we prove that
\[
I(A_1\cap\cdots\cap A_n)
\geq\sum_{j=1}^{n}I(A_j)-n+1.
\]
\begin{enumerate}[(i)]
    \item If all the events occur, the left-hand side is $1$ and the
    right-hand side is $n-n+1=1$.
    \item If at least one event does not occur, the left-hand side is $0$
    and $\sum_{j=1}^{n}I(A_j)\leq n-1$. Thus, the right-hand side is at
    most $0$, so the inequality also holds.
\end{enumerate}
Taking expectations on both sides gives
\begin{align*}
\E[I(A_1\cap\cdots\cap A_n)]
&\geq\E\left[\sum_{j=1}^{n}I(A_j)-n+1\right]\\
&=\sum_{j=1}^{n}\E[I(A_j)]-n+1.
\end{align*}
Since $\E[I(A)]=P(A)$, we obtain
\[
P(A_1\cap\cdots\cap A_n)
\geq\sum_{j=1}^{n}P(A_j)-n+1.
\]

\end{homeworkProblem}

\end{document}
